\documentclass[11pt,a4paper]{article}

%---------------------------------------------------------
%  Packages
%---------------------------------------------------------
\usepackage[utf8]{inputenc}
\usepackage[T1]{fontenc}
\usepackage{geometry}      % Margins
\geometry{margin=1in}
\usepackage[T1]{fontenc}
\usepackage{lmodern}
\usepackage{microtype}
\usepackage{hyperref}
\usepackage{enumitem}
\usepackage{amsmath,amssymb}

% Hyperref setup (optional tweaks)
\hypersetup{
    colorlinks=true,
    linkcolor=blue,
    citecolor=blue,
    urlcolor=blue
}

%---------------------------------------------------------
%  Document
%---------------------------------------------------------
\begin{document}

\section*{Key Evidence Base for the Virtual Biopsy Project}

\subsection*{Paper~1: \textit{Chen Y} \textit{et~al.}, \textit{European Radiology}, 2025\\
(Systematic Review \& Meta-Analysis)}

This comprehensive meta-analysis synthesises the results of \textbf{89 studies} on AI models that predict \textbf{EGFR, ALK, and KRAS} mutations in NSCLC.\cite{Chen2025}
The pooled performance establishes a robust benchmark:

\begin{itemize}[itemsep=2pt]
  \item Radiomics models achieved a mean sensitivity of $\sim 0.75$ for all three biomarkers.
  \item False--positive rates were non-trivial, ranging from $0.225$ (ALK) to $0.376$ (KRAS).
  \item A multimodal strategy (radiomics $+$ clinical data) boosted EGFR sensitivity to $0.800$.\cite{Chen2025b}
\end{itemize}

\textbf{Significance.} 
The study provides strong, aggregate evidence validating the \emph{Virtual Biopsy} project’s core hypothesis: multimodal architectures outperform imaging-only approaches. 
It also offers clear, quantitative benchmarks that the project must at least match—ideally surpass.

%---------------------------------------------------------
\subsection*{Multimodal Radiomics and Deep-Learning Fusion}
%---------------------------------------------------------

\subsubsection*{Paper~2: \textit{Huang L} \textit{et~al.}, \textit{Academic Radiology}, 2024\\
(Fusion of Radiomics, DL, and Clinical Features)}

Using a cohort of \textbf{826} lung adenocarcinoma patients, this study proposes an advanced fusion framework for EGFR prediction.\cite{Huang2024}
Two main innovations drive its success:

\begin{enumerate}[itemsep=2pt]
  \item \textbf{Extended Region of Interest.}  
        Image analysis encompasses the \emph{peritumoural} zone (up to 10\,mm beyond the tumour border), capturing micro-environmental signals.
  \item \textbf{Tri-modal Feature Fusion.}  
        Hand-crafted radiomic features, CNN-derived deep features, and clinical variables are stacked in a blended model.
\end{enumerate}

The final ensemble delivered an \textbf{AUC of 0.925} (internal validation) and \textbf{0.889} (external validation), significantly outperforming single-modality baselines.\cite{Huang2024}

\textbf{Significance.}  
The work supplies a technical blueprint for the Virtual Biopsy pipeline. 
Importantly, it shows that including the tumour \emph{micro-environment} can unlock additional predictive power—an advantage over physical biopsy, which samples only a small core of the lesion.

%---------------------------------------------------------
%  References  (swap for BibTeX if preferred)
%---------------------------------------------------------
\begin{thebibliography}{9}
\bibitem{Chen2025}
Chen~Y, \textit{et~al}.
Comprehensive AI prediction of driver mutations in NSCLC: a meta-analysis.
\textit{European Radiology}. 2025.

\bibitem{Chen2025b}
Chen~Y, \textit{et~al}.
Supplementary data to the above study. 2025.

\bibitem{Huang2024}
Huang~L, \textit{et~al}.
Multimodal fusion for EGFR status prediction.
\textit{Academic Radiology}. 2024.
\end{thebibliography}

\end{document}